(1) 固定端满足 \(y(0,t)=y(L,t)=0\)，波动方程 \(y_{tt}=c^2y_{xx}\)，\(c=\sqrt{T/\mu}\)。分离变量得到正交本征函数 \(\sin(n\pi x/L)\) 和频率 \(\omega_n=n\pi c/L\)。因初始速度处处为零，各模只取余弦时间因子；斜率在拨弦点不连续不妨碍傅里叶级数按能量意义收敛。
(2) 初始形状为 \(f(x)=Hx/a\)（\(x\le a\)）与 \(f(x)=H(L-x)/(L-a)\)（\(x\ge a\)）。按正弦正交关系 \(b_n=(2/L)\int_0^L f(x)\sin(n\pi x/L)\,dx\)，对两段分部积分后边界项抵消，得到 \(b_n=2HL^2\sin(n\pi a/L)/[\pi^2n^2a(L-a)]\)。系数衰减为 \(n^{-2}\)，反映初始位移连续但导数有拐角。
(3) 通解为 \(y(x,t)=\sum_{n\ge1}b_n\sin(n\pi x/L)\cos(\omega_nt)\)，并自动满足零初速度。初始弹性应变能 \(E_0=(T/2)\int_0^L[f'(x)]^2dx=(TH^2/2)(1/a+1/(L-a))\)。若 \(a/L=p/q\) 为最简分数，则当 \(n\) 为 \(q\) 的整数倍时 \(\sin(n\pi a/L)=0\)，这些模的空间结点恰落在拨弦处，无法由该初始形状激发。
(4) 第 \(n\) 模的最大应变能为 \(E_n=(T/2)b_n^2(n\pi/L)^2\int_0^L\cos^2(n\pi x/L)\,dx=Tn^2\pi^2 b_n^2/(4L)\)。正交性令不同模之间的交叉能量积分为零，故 \(\sum_nE_n=E_0\)，时间演化只在模内交换动能与势能。取 \(a=L/2\) 时 \(b_n=8H\sin(n\pi/2)/(\pi^2n^2)\)，所以全部偶数模严格为零。
